% Ruled-paper template for typesetting a scanned maths workbook.
% Copy this file for each section, fill in the title and the problems,
% and compile twice with xelatex. See SKILL.md for the full workflow.
%
% Settled format (measure the user's actual paper first; these numbers
% are for 204x291mm ruled paper with 8mm line spacing):
%   PDF page = paper size, no bleed. Text width 115mm (left 35mm,
%   right edge 150mm). Right 50mm is a correction zone, separated by a
%   thin dashed line at x=154mm, drawn on every page. No page numbers.
%   Baseline grid 8mm, first baseline 19mm from the top of the page.
%   Inline fractions use \tfrac. Problem numbers use \textcircled.
%   Answer space is always a whole multiple of 8mm.
%
% Compile: xelatex twice (the TikZ dashed line needs two passes).
% Requires texlive-xetex + texlive-lang-chinese and Noto Serif CJK SC.
\documentclass[11pt]{article}
\usepackage[paperwidth=204mm,paperheight=291mm,left=35mm,right=54mm,top=11mm,bottom=12mm]{geometry}
\usepackage{xeCJK}
\setCJKmainfont{Noto Serif CJK SC}
\usepackage{amsmath,amssymb,graphicx,xcolor}
\usepackage{tikz}
% Correction zone: thin dashed line at x=154mm (50mm from the right
% edge), running the full height of every page. Must compile twice;
% a single pass draws the line off the page. When you move the line,
% change BOTH endpoints or it draws slanted.
\AddToHook{shipout/background}{%
\begin{tikzpicture}[remember picture,overlay]
\draw[black,thin,dashed] ([xshift=154mm]current page.north west) -- ([xshift=154mm]current page.south west);
\end{tikzpicture}%
}

\newlength{\rulepitchlen}
\setlength{\rulepitchlen}{8mm}

% Round a block (knowledge box, figure, table) up to a whole number of
% 8mm lines, so the next text baseline lands back on a ruled line.
% Re-measure the offset for each new box height or figure size; do not
% copy a calibration value from another block.
\newsavebox{\gbbox}
\newlength{\gbh}
\newlength{\gbrem}
\newlength{\gbgap}
\newcount{\gbn}
\newcommand{\gridblock}[1]{%
\par\noindent
\sbox{\gbbox}{\parbox[t]{\textwidth}{#1}}%
\setlength{\gbh}{\dimexpr\ht\gbbox+\dp\gbbox\relax}%
\gbn=\numexpr\gbh/\rulepitchlen\relax
\setlength{\gbrem}{\dimexpr\gbh-\gbn\rulepitchlen\relax}%
\ifdim\gbrem>0pt \advance\gbn by 1 \fi
\usebox{\gbbox}\par\nointerlineskip
\setlength{\gbgap}{\dimexpr\gbn\rulepitchlen-\dp\gbbox-\topskip\relax}%
\vspace{\gbgap}%
}

\setlength{\parindent}{0pt}

\pagestyle{empty}
\begin{document}
% One 8mm line, no extra space before or after paragraphs, first
% baseline 8mm below the top of the text area (19mm from the page top).
\fontsize{11pt}{8mm}\selectfont
\setlength{\parskip}{0pt}
\setlength{\topskip}{8mm}
\setlength{\lineskip}{0pt}
\setlength{\lineskiplimit}{0pt}

% ================= 题目 =================
\noindent{\fontsize{16pt}{8mm}\selectfont\bfseries 第1章\ 有理数\par}
\noindent{\fontsize{14pt}{8mm}\selectfont\bfseries 第一周\ 正数和负数\ 数轴\ 相反数\par}
\noindent{\fontsize{12pt}{8mm}\selectfont\bfseries 一、选择题\par}
\noindent\textcircled{1} 在数轴上，原点及原点右边的点所表示的数是（\hspace{1.5em}）。

\noindent （A）正数 \hfill （B）负数 \hfill （C）非正数 \hfill （D）非负数

% Choice / fill-in problem: leave 1 line (8mm) after it.
\vspace{8mm}

\noindent\textcircled{2} 已知 $a$、$b$ 在数轴上的对应点的位置如图所示，那么下列各式中错误的是（\hspace{1.5em}）。

% Figure in a \gridblock so the text after it stays on the grid.
% Crop the figure from a 300 DPI render of the source page; enlarge
% only, never redraw. Add a centred caption naming the figure.
\gridblock{%
{\centering\includegraphics[width=6.1cm]{figures/fig2.png}\par}
{\centering\small 第 2 题图\par}%
}

\noindent （A）$a<0$ \hfill （B）$b>0$ \hfill （C）$-a>0$ \hfill （D）$-b>0$

% Choice / fill-in problem: leave 1 line (8mm) after it.
\vspace{8mm}

\noindent{\fontsize{12pt}{8mm}\selectfont\bfseries 二、填空题\par}
\noindent\textcircled{3} 如果 $a+3$ 的相反数是 $-3$，那么 $a=$\underline{\hspace{3em}}。

% Choice / fill-in problem: leave 1 line (8mm) after it.
\vspace{8mm}

\noindent{\fontsize{12pt}{8mm}\selectfont\bfseries 三、解答题\par}
\noindent\textcircled{4} 把 $-1$、$2.6$、$-\tfrac{22}{7}$、$0.35$、$0$、$25$、$0.\dot{4}\dot{4}$ 填在相应的横线上。

\noindent 正有理数：\underline{\hspace{20em}}；

\noindent 非负整数：\underline{\hspace{20em}}；

\noindent 非正整数：\underline{\hspace{20em}}。

% Worked problem: leave 2 lines (16mm) after it for working.
\vspace{16mm}

\end{document}
